\section{Intervjuförberedelse}
\label{sec:intervju}

Tekniska intervjufrågor på en investmentbank är förutsägbara. Det som skiljer kandidater är
inte om de kan svaret utan om de kan ge det strukturerat på en minut. Varje svar nedan är
skrivet för att kunna sägas högt på 60 till 90 sekunder: först rubriken, sedan stegen, sist en
nyans som visar att du förstår. Frågorna står på engelska, som de ställs. Svaren står på
svenska så att du lär dig innehållet och inte en ramsa, följda av de engelska fraser du ska
kunna säga utan att tveka.

\begin{intervjutips}
Öva högt och med klocka. Börja varje svar med strukturen (”There are three steps”) och sluta
när du är klar. Om du inte kan en sak: säg vad du vet och resonera dig fram. Gissa aldrig en
siffra som om den vore säker.
\end{intervjutips}

\subsection{Redovisning}

\begin{fraga}{Walk me through the three financial statements.}
Resultaträkningen visar intäkter minus kostnader under en period och slutar i nettoresultatet.
Balansräkningen visar tillgångar, skulder och eget kapital en viss dag, och tillgångarna är
alltid lika med skulder plus eget kapital. Kassaflödesanalysen börjar med nettoresultatet,
justerar för poster som inte är betalningar och för förändringar i rörelsekapital, lägger till
investeringar och finansiering och slutar i periodens förändring av kassan. Kopplingarna:
nettoresultatet går in i kassaflödet och i eget kapital, och kassaflödets slutrad blir kassan i
balansräkningen.
\fraser{Net income flows into retained earnings and is the starting point of the cash flow
statement. The ending cash balance ties back to the balance sheet.}
\end{fraga}

\begin{fraga}{How does a 10 increase in depreciation affect the three statements?}
Resultaträkningen: resultatet före skatt faller med 10. Med 20,6\pct{} skatt faller skatten med
2,06 och nettoresultatet med 7,94. Kassaflödet: nettoresultatet är 7,94 lägre, men
avskrivningen på 10 läggs tillbaka eftersom den inte är en utbetalning, så kassaflödet från
rörelsen och kassan ökar med 2,06. Balansräkningen: anläggningstillgångarna minskar med 10 och
kassan ökar med 2,06, alltså tillgångar $-7,94$. Eget kapital minskar med 7,94 genom det lägre
resultatet. Det går ihop. Se tabell~\ref{tab:avskrivning}.
\fraser{Depreciation is a non-cash expense, so cash actually goes up by the tax shield. Both
sides of the balance sheet are down by 7.94, so it balances.}
\end{fraga}

\begin{fraga}{What is EBITDA, and why do bankers use it?}
EBITDA är rörelseresultatet före avskrivningar, räntor och skatt. Det används som en grov
approximation av rörelsens kassaflöde och för att jämföra bolag, eftersom det inte påverkas av
skuldsättning, skattesituation eller avskrivningsprinciper. Svagheten är att det bortser från
investeringar, rörelsekapital och skatt. Ett kapitalintensivt bolag kan ha hög EBITDA och ändå
lågt fritt kassaflöde. Nyansen: sedan IFRS~16 ligger leasingkostnader under EBITDA, vilket
lyfter måttet för bolag med mycket leasing.
\fraser{It is a proxy for operating cash flow that is independent of capital structure. But it
ignores capex, working capital and taxes.}
\end{fraga}

\begin{fraga}{What is the difference between unlevered and levered free cash flow?}
Det obelånade kassaflödet, FCFF, tillhör alla finansiärer: EBIT efter skatt plus avskrivningar
minus capex minus ökning av rörelsekapital. Det belånade, FCFE, är det som blir kvar till
aktieägarna efter räntor och nettoupplåning. FCFF diskonteras med WACC och ger rörelsevärdet,
FCFE med avkastningskravet på eget kapital och ger aktievärdet. I en vanlig DCF används FCFF,
eftersom man då slipper prognostisera skuldsättningen.
\fraser{Unlevered free cash flow is available to all capital providers and is discounted at the
WACC. Levered free cash flow belongs to equity holders and is discounted at the cost of
equity.}
\end{fraga}

\begin{fraga}{What is working capital, and how does a change in it affect cash flow?}
Rörelsekapital är omsättningstillgångar minus kortfristiga rörelseskulder, i sin enklaste form
kundfordringar plus lager minus leverantörsskulder. En ökning binder pengar och är ett
kassautflöde: om lagret ökar med 10 har bolaget betalat 10 som ännu inte syns som kostnad.
Växande bolag behöver därför normalt finansiera sitt rörelsekapital. Bolag som får betalt i
förskott kan ha negativt rörelsekapital, och då frigör tillväxten pengar.
\fraser{An increase in net working capital is a use of cash.}
\end{fraga}

\subsection{DCF och WACC}

\begin{fraga}{Walk me through a DCF.}
En DCF värderar ett bolag som nuvärdet av dess framtida kassaflöden. Fem steg. Ett: prognostisera
det obelånade fria kassaflödet i fem till tio år. Två: beräkna WACC. Tre: beräkna ett slutvärde
för tiden efter prognosen, med Gordons formel eller en exit-multipel. Fyra: diskontera
kassaflödena och slutvärdet med WACC och summera till rörelsevärdet. Fem: dra av nettoskuld och
minoritetsintressen för att få aktievärdet, och dela med antalet aktier. Avsluta med en
känslighetsanalys på WACC och tillväxt, eftersom slutvärdet oftast är det mesta av värdet.
\fraser{I project unlevered free cash flows, discount them at the WACC, add a terminal value,
and bridge from enterprise value to equity value.}
\end{fraga}

\begin{fraga}{What is WACC and how do you calculate it?}
WACC är det vägda avkastningskravet från alla som finansierar bolaget och används som
diskonteringsränta för FCFF. Det är andelen eget kapital gånger avkastningskravet på eget
kapital, plus andelen skuld gånger låneräntan efter skatt, se~\eqref{eq:wacc}. Vikterna är
marknadsvärden. Avkastningskravet på eget kapital kommer från CAPM. Låneräntan multipliceras
med ett minus skattesatsen eftersom räntan är avdragsgill. För Atlas Copco fick jag 7,65\pct:
96\pct{} eget kapital till 7,84\pct{} och 4\pct{} skuld till 2,84\pct{} efter skatt.
\fraser{Cost of equity times the equity weight, plus after-tax cost of debt times the debt
weight, using market values.}
\end{fraga}

\begin{fraga}{How do you calculate the cost of equity?}
Med CAPM: riskfri ränta plus beta gånger marknadens riskpremie. Den riskfria räntan är en lång
statsobligation i samma valuta som kassaflödena, här den svenska tioåringen på 3,075\pct.
Riskpremien är vad aktiemarknaden väntas ge utöver det. Jag använde Damodarans 4,23\pct{} för
Sverige. Beta mäter aktiens känslighet för marknaden. Tanken är att investerare bara får
betalt för risk som inte kan diversifieras bort. Nyansen: riskpremien går inte att observera,
och ett annat källval ger ett annat värde.
\fraser{Risk-free rate plus beta times the equity risk premium. Only systematic risk is
compensated.}
\end{fraga}

\begin{fraga}{What is beta, and how do you unlever and relever it?}
Beta är lutningen i en regression av aktiens avkastning på marknadens, kovariansen delat med
marknadens varians. Ett observerat beta speglar både rörelsens risk och skuldsättningen. För
att låna beta från jämförelsebolag rensar man därför först bort deras skuldsättning: beta delat
med ett plus ett minus skattesatsen gånger skuld genom eget kapital. Man tar medianen och
lägger sedan på målbolagets egen kapitalstruktur med samma formel baklänges,
se~\eqref{eq:hamada}. Skälet att använda jämförelsebolag är att ett enskilt regressionsbeta är
brusigt.
\fraser{Unlevered beta equals levered beta divided by one plus one minus the tax rate times
debt to equity. Then I relever at the target's capital structure.}
\end{fraga}

\begin{fraga}{How do you calculate terminal value, and which method do you prefer?}
Två metoder. Gordons formel antar evig tillväxt: nästa års kassaflöde delat med WACC minus
tillväxt. Exit-multipeln antar att bolaget säljs sista prognosåret till en multipel av EBITDA.
Jag räknar alltid båda och stämmer av dem mot varandra: vilken multipel motsvarar min tillväxt,
och vilken tillväxt motsvarar min multipel? I mitt case gav 2\pct{} tillväxt en multipel på
11,7 gånger, medan jämförelsebolagens 19,4 gånger motsvarade 4,3\pct{} evig tillväxt. Den
skillnaden var hela förklaringen till att metoderna inte var överens.
\fraser{The perpetuity growth method and the exit multiple method. I cross-check one against the
other through the implied multiple and the implied growth rate.}
\end{fraga}

\begin{fraga}{What is the mid-year convention?}
En vanlig DCF diskonterar varje års kassaflöde som om det kom vid årets slut. I verkligheten
kommer det löpande under året, i snitt vid halvårsskiftet. Halvårskonventionen diskonterar
därför år $t$ med $t$ minus ett halvt år, vilket höjer värdet något. Det viktiga är
slutvärdet: Gordon-värdet ska också justeras upp med ett halvt års ränta, medan ett slutvärde
från en exit-multipel avser en försäljning vid årets slut och inte ska justeras.
\fraser{Cash flows arrive evenly through the year, so I discount them at t minus one half. The
exit multiple terminal value gets no mid-year adjustment.}
\end{fraga}

\begin{fraga}{What is the difference between enterprise value and equity value?}
Rörelsevärdet är värdet av hela rörelsen för alla finansiärer, oavsett hur den är finansierad.
Aktievärdet är det som tillhör aktieägarna. Bryggan: rörelsevärde minus nettoskuld, minus
minoritetsintressen och preferensaktier, är aktievärde. Åt andra hållet: börsvärde plus
nettoskuld är rörelsevärde. Rörelsevärdet används med mått före räntor, som EBITDA, och
aktievärdet med mått efter räntor, som nettoresultat.
\fraser{Enterprise value is the value of the operations to all investors. Equity value is what
is left for shareholders after net debt and other claims.}
\end{fraga}

\begin{fraga}{Why do you subtract cash and add minority interest when you calculate enterprise value?}
Kassan dras av eftersom den inte behövs i rörelsen: en köpare som betalar för aktierna får
kassan med på köpet och kan använda den till att betala skulden. Minoritetsintressen läggs
till eftersom koncernens EBITDA räknas till 100\pct{} även för dotterbolag som delvis ägs av
andra. Täljaren måste då också innehålla hela värdet, annars jämför man äpplen med päron.
\fraser{Cash is a non-operating asset. Minority interest is added for consistency, because
consolidated EBITDA includes one hundred percent of the subsidiary.}
\end{fraga}

\begin{fraga}{Which has a bigger impact on a DCF: the discount rate or the growth rate?}
Båda verkar genom samma nämnare, WACC minus tillväxt, så en procentenhet i den ena är ungefär
lika mycket värd som en procentenhet i den andra för slutvärdet. Men WACC påverkar dessutom
diskonteringen av alla kassaflöden, så WACC väger normalt något tyngre. Det avgörande är hur
liten nämnaren är. I mitt case var den 5,65 procentenheter, och då sänker en procentenhet högre
WACC slutvärdet med omkring 15\pct.
\fraser{The value is driven by the spread between the WACC and the growth rate. The smaller the
spread, the more sensitive the valuation.}
\end{fraga}

\begin{fraga}{Terminal value is 80 percent of your DCF. Is that a problem?}
Det är normalt för ett moget bolag med fem års prognos, och i mitt case var andelen 76 till
84\pct. Men det betyder att värderingen står och faller med två eller tre antaganden om tiden
efter år fem. Jag hanterar det på tre sätt: stämmer av de två slutvärdesmetoderna mot varandra,
visar en känslighetstabell, och räknar baklänges vad marknadspriset kräver. En längre
prognosperiod minskar andelen men flyttar bara osäkerheten.
\fraser{That is typical, and it is exactly why I sanity-check the implied multiple and run
sensitivities.}
\end{fraga}

\begin{fraga}{Does more debt always lower the WACC?}
Nej. Skuld är billigare än eget kapital och räntan är avdragsgill, så lite skuld sänker WACC.
Men mer skuld gör aktien mer riskfylld: beta stiger enligt Hamada, och därmed avkastningskravet
på eget kapital. Vid hög skuldsättning stiger också låneräntan och risken för finansiella
problem. WACC är därför U-formad i skuldsättningen. Utan skatter och konkurskostnader tar
effekterna ut varandra helt. Det är Modigliani och Millers resultat.
\fraser{Debt is cheaper and tax-deductible, but it raises the cost of equity and eventually the
cost of debt. There is an optimal capital structure.}
\end{fraga}

\subsection{Multiplar och metodval}

\begin{fraga}{What are the main valuation methods, and which gives the highest value?}
Tre huvudmetoder: DCF, som ger ett fristående värde ur kassaflödena, jämförelsebolag, som visar
vad börsen betalar för liknande bolag, och jämförbara transaktioner \en{precedent transactions},
som visar vad köpare har betalat för hela bolag. Därtill LBO-analys. Som tumregel ger
transaktioner högst värde, eftersom priserna innehåller en kontrollpremie, och LBO lägst,
eftersom en finansiell köpare kräver hög avkastning. DCF kan hamna var som helst beroende på
antagandena. Mitt projekt har inte transaktionsmultiplar, eftersom det saknas gratis data.
\fraser{DCF, trading comps and precedent transactions. Precedents usually come out highest
because they include a control premium.}
\end{fraga}

\begin{fraga}{When would you use EV/EBITDA rather than P/E?}
EV/EBITDA när jag vill jämföra rörelser oberoende av finansiering: måttet påverkas inte av
skuldsättning, skatt eller avskrivningsprinciper, och det är standard för industribolag och i
transaktioner. P/E när kapitalstrukturen är en del av verksamheten, framför allt banker och
försäkringsbolag, där rörelsevärde och EBITDA inte är meningsfulla. P/E är också det mått
aktieinvesterare oftast tittar på. Viktigast är att inte blanda: rörelsevärde hör ihop med
resultat före räntor, börsvärde med resultat efter räntor.
\fraser{EV/EBITDA is capital structure neutral. P/E is an equity multiple and is affected by
leverage.}
\end{fraga}

\begin{fraga}{How do you select comparable companies?}
Jag letar efter bolag med samma affärsmodell och samma drivkrafter: bransch och geografi först,
sedan storlek, tillväxt, marginaler och risk. Hellre fem bra jämförelsebolag än femton
dåliga. Sedan rensar jag uteliggare och tittar på median och kvartiler, inte medelvärde. I mitt
projekt gjorde en algoritm urvalet inom sektorn, och det lärde mig var gränsen går: den
hittade rimliga bolag men kan inte veta att två bolag med samma nyckeltal säljer helt olika
saker.
\fraser{Same industry and business model first, then size, growth, margins and risk. I use the
median rather than the mean.}
\end{fraga}

\begin{fraga}{Why would two similar companies trade at different multiples?}
En multipel är en hoptryckt DCF: den bestäms av tillväxt, risk och hur mycket av resultatet som
blir kassaflöde, se~\eqref{eq:multipel-dcf}. Högre väntad tillväxt, lägre risk eller högre
avkastning på kapitalet motiverar en högre multipel. Därutöver kan skillnaden vara teknisk:
olika redovisning, engångsposter i resultatet eller olika perioder. Först när allt det är
uteslutet kan man säga att ett av bolagen är felprissatt.
\fraser{Differences in growth, risk and returns on capital. A multiple is shorthand for a DCF.}
\end{fraga}

\begin{fraga}{A company has negative earnings. How do you value it?}
P/E fungerar inte, och är EBITDA också negativ faller EV/EBITDA bort. Då återstår
omsättningsmultiplar, med reservationen att de bortser från lönsamheten, och en DCF med en
uttalad väg till lönsamhet. En sådan DCF är mycket känslig för när och på vilken nivå
marginalen stabiliseras. Min motor avstår helt från DCF vid negativ EBITDA, eftersom den bara
kan skriva fram historiken, och det gällde 154 av 687 svenska bolag. Hellre det än ett tal
utan mening.
\fraser{I would look at revenue multiples and a DCF with an explicit path to profitability, and
be clear about the uncertainty.}
\end{fraga}

\begin{fraga}{How would you value a bank?}
Inte med en vanlig DCF. För en bank är skulden råvaran, inte finansieringen: inlåning och
upplåning är en del av rörelsen, så rörelsevärde, EBITDA och FCFF saknar mening. Man värderar
i stället det egna kapitalet direkt, med P/E och kurs mot eget kapital (P/B) ställt mot
avkastningen på eget kapital, eller med en utdelningsmodell där utdelningarna diskonteras med
avkastningskravet på eget kapital. Min motor ger banker bara P/E, och säger det öppet.
\fraser{For a bank, debt is a raw material rather than a source of financing, so I use equity
multiples like P/E and price to book, or a dividend discount model.}
\end{fraga}

\begin{fraga}{What is a football field chart?}
Ett diagram med en liggande stapel per värderingsmetod som visar metodens intervall för värde
per aktie, och en lodrät linje för aktuell kurs eller ett bud. Det används i pitchar och i
underlag till styrelser för att visa var metoderna överlappar. Poängen är inte ett exakt tal
utan ett försvarbart intervall. Man ska kunna förklara varje stapels ändpunkter och varför
staplarna ligger där de gör.
\fraser{It summarises the valuation ranges from each methodology side by side, against the
current share price.}
\end{fraga}

\subsection{LBO}

\begin{fraga}{Walk me through an LBO.}
I en LBO köper ett riskkapitalbolag ett bolag med en stor andel lånade pengar. Stegen: ett,
anta ett köppris, oftast som multipel av EBITDA. Två, bestäm finansieringen i en uppställning
av källor och användning: lån i olika skikt och resten eget kapital. Tre, prognostisera
kassaflödet och använd det till ränta och amortering. Fyra, anta en försäljning efter omkring
fem år till en exit-multipel, dra av kvarvarande skuld och få ägarens likvid. Fem, räkna IRR
och MOIC på ägarens insats. Ofta vänder man på det och löser ut det pris som ger ett visst
avkastningskrav.
\fraser{Purchase price, sources and uses, debt paydown from free cash flow, exit, and then the
sponsor's IRR and multiple of money.}
\end{fraga}

\begin{fraga}{Why does an LBO use debt?}
Hävstång. Skuld är billigare än eget kapital, och om bolaget är värt mer än skulden vid
försäljningen tillfaller hela värdeökningen ägaren, som bara har satsat en del av köpeskillingen.
Samma värdeökning på en mindre insats ger högre avkastning. Dessutom är räntan avdragsgill, och
kravet att betala ränta tvingar fram disciplin i bolaget. Baksidan är risken: skulden ska
betalas även om resultatet faller, och då är det ägarens kapital som försvinner först.
\fraser{Leverage amplifies equity returns, and interest is tax-deductible. The trade-off is
financial risk.}
\end{fraga}

\begin{fraga}{What drives returns in an LBO?}
Tre saker. Resultattillväxt: bolaget säljs med högre EBITDA än det köptes med. Multipelexpansion:
det säljs till en högre multipel än köpet. Amortering: bolagets kassaflöde betalar ned skulden,
så en större del av rörelsevärdet tillhör ägaren. De tre summerar exakt till värdeskapandet
efter avgifter, se~\eqref{eq:attribution}. Den mest pålitliga källan är amorteringen och den
minst pålitliga multipeln. I mitt case köptes bolaget till 28 gånger EBITDA och såldes till 19,
och multipelfallet åt upp allt annat.
\fraser{EBITDA growth, multiple expansion and debt paydown.}
\end{fraga}

\begin{fraga}{What makes a good LBO candidate?}
Ett bolag som klarar skuld: stabila och förutsägbara kassaflöden, låga investeringsbehov, lågt
rörelsekapitalbehov, stark marknadsposition och låg befintlig skuldsättning. Det ska gå att
köpa till en rimlig multipel, det ska finnas förbättringsmöjligheter i rörelsen, och det ska
finnas en trolig väg ut om några år, till en industriell köpare, en annan fond eller börsen.
Cykliska bolag och bolag med stora investeringsbehov passar sämre.
\fraser{Stable cash flows, low capex requirements, a strong market position and a clear exit
route.}
\end{fraga}

\begin{fraga}{What is the difference between IRR and MOIC?}
MOIC är pengar ut delat med pengar in och bryr sig inte om tiden. IRR är den årliga avkastning
som gör nuvärdet av alla ägarens kassaflöden till noll. Med en insats och en utbetalning är
IRR lika med MOIC upphöjt till ett genom antalet år, minus ett. Dubbla pengarna är 26\pct{} om
året på tre år men 15\pct{} på fem. En tidig utdelning höjer IRR utan att höja MOIC, och därför
tittar investerare på båda.
\fraser{MOIC tells you how much, IRR tells you how fast. Two times your money in five years is
roughly a fifteen percent IRR.}
\end{fraga}

\begin{fraga}{What is the difference between senior and subordinated debt, and what is a cash sweep?}
Seniora lån har förmånsrätt: de betalas först vid en konkurs och har ofta säkerhet i bolagets
tillgångar. De har därför lägre ränta, amorteras löpande och har hårdare villkor. Efterställda
lån står längre bak i kön, har högre ränta och betalas ofta i en klump vid förfall. Ett
\en{cash sweep} betyder att överskottskassa, utöver den obligatoriska amorteringen, används
till att betala det seniora lånet i förtid. Det sänker räntekostnaden och risken, och det
ökar ägarens andel av värdet.
\fraser{Senior debt ranks first and is cheaper. A cash sweep uses excess cash flow to prepay
senior debt.}
\end{fraga}

\subsection{Omdöme}

\begin{fraga}{Your DCF says 116 and the stock trades at 192. Who is right?}
Det kan modellen inte avgöra, och jag skulle inte påstå det. Det den kan säga är vad kursen
förutsätter. Jag räknade baklänges: kursen kräver ungefär 1,4 procentenheter lägre
avkastningskrav och en procentenhet högre evig tillväxt än jag antog. Det är inte orimligt för
ett kvalitetsbolag, och min modell är snål: den bromsar tillväxten till 2\pct{} redan år fem.
Samtidigt såg jag samma mönster i hela universumet, där DCF:en värderar storbolag till 0,76
av kursen i median. Det tyder på att gapet säger mer om modellens antaganden än om just det
här bolaget.
\fraser{A DCF does not tell you the price is wrong. It tells you what you need to believe for
the price to be right.}
\end{fraga}

\subsection{Att presentera projektet på en minut}

Du kommer att få frågan ”tell me about something you have built” eller ”walk me through a
project on your CV”. Svaret ska vara på engelska, ta en minut och sluta i en insikt, inte i en
lista över tekniker.

\begin{tcolorbox}[enhanced,colback=accent!4,colframe=accent,boxrule=0.5pt,arc=2pt,
  left=10pt,right=10pt,top=8pt,bottom=8pt,title={The pitch (about 60 seconds)},
  fonttitle=\sffamily\bfseries\small,coltitle=white,colbacktitle=accent,
  before upper={\setlength{\parskip}{0.5\baselineskip}}]
\begin{english}
I built a valuation engine in Python that values every
Swedish listed company on Yahoo Finance, 687 of them, the way a bank would: a DCF with a
CAPM-based WACC, trading comps, an LBO and a Monte Carlo simulation, summarised in a football
field and exported to Excel.

I then used it for a case study on Atlas Copco. The interesting result was not a price target
but a disagreement. The DCF with two percent terminal growth gives about 116 kronor per share.
Comps and the exit-multiple DCF give around 175. The stock trades at 192. I traced the gap to
the terminal value: the peer multiple of 19 times EBITDA implies more than four percent growth
forever. So the model cannot tell you who is right, but it shows exactly what the market price
assumes.

What I learned most from was where the mechanical approach breaks. For example, the clustering
barely selects peers on its own, and the Monte Carlo range mostly reflects the volatilities I
assumed.
\end{english}
\end{tcolorbox}

\paragraph{Varför den fungerar.} Den börjar med vad projektet \emph{är} i en mening, med
siffror som visar omfånget. Den går sedan till ett resultat som kräver förståelse (varför
metoderna är oense), inte till en uppräkning av bibliotek. Och den slutar med svagheter som du
själv har hittat. Det sista är det viktigaste: den som intervjuar har sett många DCF-modeller,
men få kandidater som självmant säger var deras modell är svag.

\paragraph{Följdfrågor du ska räkna med, och ärliga svar.}

\begin{description}
  \item[”Would you buy the stock?”] Nej, inte på det här underlaget. Modellen är ett läroexempel
    på gratis data som inte är avstämda mot årsredovisningen. Den visar vad kursen kräver,
    inte om kursen är fel.
  \item[”Why k-means? Did it work?”] Bara delvis. Den valde två kluster, och 119 av 148
    kandidater hamnade i samma som Atlas Copco. Urvalet görs i praktiken av närmaste grannar
    inom sektorn med förtur för samma bransch. En analytiker hade valt för hand, och ABB hade
    troligen kommit med.
  \item[”Where does your size premium come from?”] Det är ett eget antagande, inte skattat ur
    data och inte källbelagt. Det finns där för att regressionsbeta för illikvida småbolag blir
    för låg. Det går att stänga av, och i caset är premien noll.
  \item[”What does the Monte Carlo actually tell you?”] Hur känslig modellen är för sina egna
    indata. Bredden bestäms av standardavvikelser som jag har antagit, och tillväxtens bygger
    på tre observationer. Att 3,3\pct{} av dragningarna ligger över kursen är ingen sannolikhet
    för att aktien är undervärderad.
  \item[”What would you do next?”] Rullande tolvmånaderssiffror i stället för senaste bokslut,
    en stubbperiod, justering för leasing, kurs mot eget kapital för banker och fastighetsbolag,
    transaktionsmultiplar, och framför allt en avstämning av indata mot bolagens egna rapporter.
\end{description}
